\documentclass[10pt,a4paper]{amsart}
\usepackage{amsmath,amssymb,amsthm,graphicx}
% Retained preamble for the Stacks Project chapter snapshot in this directory.
% The chapter's own preamble.tex is a copy of the shared Stacks Project preamble; the block
% below is sliced MECHANICALLY, line for line, from it. Nothing was added, removed or
% redefined; the wrapper class, the xr-hyper/hyperref wiring and the externaldocument list
% are omitted because the assembler supplies the document class.
%
% lines taken from preamble.tex: 3, 5, 9, 11, 12, 156, 157, 158, ... (51 lines total)
%
% package coverage restored from the chapter's own preamble (2026-10-13): xy, xy-2cell

\usepackage[all]{xy}
\xyoption{2cell}
\UseAllTwocells

\newenvironment{reference}{\comment}{\endcomment}
\newenvironment{slogan}{\comment}{\endcomment}
\newenvironment{history}{\comment}{\endcomment}
\theoremstyle{plain}
\newtheorem{theorem}[subsection]{Theorem}
\newtheorem{proposition}[subsection]{Proposition}
\newtheorem{lemma}[subsection]{Lemma}
\theoremstyle{definition}
\newtheorem{definition}[subsection]{Definition}
\newtheorem{example}[subsection]{Example}
\newtheorem{exercise}[subsection]{Exercise}
\newtheorem{situation}[subsection]{Situation}
\theoremstyle{remark}
\newtheorem{remark}[subsection]{Remark}
\newtheorem{remarks}[subsection]{Remarks}
\numberwithin{equation}{subsection}
\def\lim{\mathop{\mathrm{lim}}\nolimits}
\def\colim{\mathop{\mathrm{colim}}\nolimits}
\def\Spec{\mathop{\mathrm{Spec}}}
\def\Hom{\mathop{\mathrm{Hom}}\nolimits}
\def\Ext{\mathop{\mathrm{Ext}}\nolimits}
\def\SheafHom{\mathop{\mathcal{H}\!\mathit{om}}\nolimits}
\def\SheafExt{\mathop{\mathcal{E}\!\mathit{xt}}\nolimits}
\def\Sch{\mathit{Sch}}
\def\Mor{\mathop{\mathrm{Mor}}\nolimits}
\def\Ob{\mathop{\mathrm{Ob}}\nolimits}
\def\Sh{\mathop{\mathit{Sh}}\nolimits}
\def\NL{\mathop{N\!L}\nolimits}
\def\CH{\mathop{\mathrm{CH}}\nolimits}
\def\proetale{{pro\text{-}\acute{e}tale}}
\def\etale{{\acute{e}tale}}
\def\QCoh{\mathit{QCoh}}
\def\Ker{\mathop{\mathrm{Ker}}}
\def\Im{\mathop{\mathrm{Im}}}
\def\Coker{\mathop{\mathrm{Coker}}}
\def\Coim{\mathop{\mathrm{Coim}}}
\DeclareMathSymbol{\boxtimes}{\mathbin}{AMSa}{"02}
\def\QCohstack{\mathcal{QC}\!\mathit{oh}}
\def\Cohstack{\mathcal{C}\!\mathit{oh}}
\def\Spacesstack{\mathcal{S}\!\mathit{paces}}
\def\Quotfunctor{\mathrm{Quot}}
\def\Hilbfunctor{\mathrm{Hilb}}
\def\Curvesstack{\mathcal{C}\!\mathit{urves}}
\def\Polarizedstack{\mathcal{P}\!\mathit{olarized}}
\def\Complexesstack{\mathcal{C}\!\mathit{omplexes}}
\def\Pic{\mathop{\mathrm{Pic}}\nolimits}
\def\Picardstack{\mathcal{P}\!\mathit{ic}}
\def\Picardfunctor{\mathrm{Pic}}
\def\Deformationcategory{\mathcal{D}\!\mathit{ef}}
\usepackage{cleveref}
\crefname{originallemma}{original lemma}{original lemmas}
\Crefname{originallemma}{Original lemma}{Original lemmas}
\makeatletter
\newlabel{algebra-lemma-characterize-henselian}{{04GG}{source}{}{sourcepointer-1}{}}
\newlabel{algebra-lemma-characterize-henselian@cref}{{[originallemma][0][]04GG}{[1][0][]source}{}{}{}}
\makeatother
\crefname{originallemma}{original lemma}{original lemmas}
\Crefname{originallemma}{Original lemma}{Original lemmas}
\makeatletter
\newlabel{algebra-lemma-flat-local-given-residue-field}{{03C3}{source}{}{sourcepointer-2}{}}
\newlabel{algebra-lemma-flat-local-given-residue-field@cref}{{[originallemma][0][]03C3}{[1][0][]source}{}{}{}}
\makeatother
\crefname{originallemma}{original lemma}{original lemmas}
\Crefname{originallemma}{Original lemma}{Original lemmas}
\makeatletter
\newlabel{example-universal-finite-syntomic}{{0FKZ}{source}{}{sourcepointer-3}{}}
\newlabel{example-universal-finite-syntomic@cref}{{[originallemma][0][]0FKZ}{[1][0][]source}{}{}{}}
\makeatother
\crefname{originallemma}{original lemma}{original lemmas}
\Crefname{originallemma}{Original lemma}{Original lemmas}
\makeatletter
\newlabel{example-universal-quasi-finite-syntomic}{{0FK8}{source}{}{sourcepointer-4}{}}
\newlabel{example-universal-quasi-finite-syntomic@cref}{{[originallemma][0][]0FK8}{[1][0][]source}{}{}{}}
\makeatother
\crefname{originallemma}{original lemma}{original lemmas}
\Crefname{originallemma}{Original lemma}{Original lemmas}
\makeatletter
\newlabel{lemma-syntomic-finite}{{0FKY}{source}{}{sourcepointer-5}{}}
\newlabel{lemma-syntomic-finite@cref}{{[originallemma][0][]0FKY}{[1][0][]source}{}{}{}}
\makeatother
\crefname{originallemma}{original lemma}{original lemmas}
\Crefname{originallemma}{Original lemma}{Original lemmas}
\makeatletter
\newlabel{lemma-universal-finite-syntomic}{{0FL0}{source}{}{sourcepointer-6}{}}
\newlabel{lemma-universal-finite-syntomic@cref}{{[originallemma][0][]0FL0}{[1][0][]source}{}{}{}}
\makeatother
\crefname{originallemma}{original lemma}{original lemmas}
\Crefname{originallemma}{Original lemma}{Original lemmas}
\makeatletter
\newlabel{lemma-universal-finite-syntomic-smooth}{{0FL1}{source}{}{sourcepointer-7}{}}
\newlabel{lemma-universal-finite-syntomic-smooth@cref}{{[originallemma][0][]0FL1}{[1][0][]source}{}{}{}}
\makeatother
\crefname{originallemma}{original lemma}{original lemmas}
\Crefname{originallemma}{Original lemma}{Original lemmas}
\makeatletter
\newlabel{lemma-universal-quasi-finite-syntomic-etale}{{0FK9}{source}{}{sourcepointer-8}{}}
\newlabel{lemma-universal-quasi-finite-syntomic-etale@cref}{{[originallemma][0][]0FK9}{[1][0][]source}{}{}{}}
\makeatother
\crefname{originallemma}{original lemma}{original lemmas}
\Crefname{originallemma}{Original lemma}{Original lemmas}
\makeatletter
\newlabel{more-algebra-lemma-henselization-regular}{{06LN}{source}{}{sourcepointer-9}{}}
\newlabel{more-algebra-lemma-henselization-regular@cref}{{[originallemma][0][]06LN}{[1][0][]source}{}{}{}}
\makeatother
\crefname{originallemma}{original lemma}{original lemmas}
\Crefname{originallemma}{Original lemma}{Original lemmas}
\makeatletter
\newlabel{more-morphisms-section-etale-localization}{{04HF}{source}{}{sourcepointer-10}{}}
\newlabel{more-morphisms-section-etale-localization@cref}{{[originallemma][0][]04HF}{[1][0][]source}{}{}{}}
\makeatother
\crefname{originallemma}{original lemma}{original lemmas}
\Crefname{originallemma}{Original lemma}{Original lemmas}
\makeatletter
\newlabel{morphisms-lemma-set-points-where-fibres-etale}{{0476}{source}{}{sourcepointer-11}{}}
\newlabel{morphisms-lemma-set-points-where-fibres-etale@cref}{{[originallemma][0][]0476}{[1][0][]source}{}{}{}}
\makeatother
\title{Stacks, Discriminants: lemma-universal-finite-syntomic-etale}
\author{The Stacks Project Authors}
\date{stacks-project:discriminant.tex:local-snapshot-20261012; GNU FDL 1.2 or later}
\begin{document}
\maketitle
\noindent\textit{Editorial note.} Statement span is lines 2417-2423 of sources/stacks-discriminant/pinned-original-source.tex; proof span is lines 2425-2479, the complete author proof (53 body lines).\par\medskip
\noindent\textit{Editorial note.} Navigation stubs: 11 cited-but-not-included labels are supplied as non-load-bearing pointers with their pinned tags where the snapshot tag index has them.\par\medskip
\noindent\textit{Editorial note.} extra\_packages is empty on purpose: the retained preamble already loads xy with options and a second declaration clashes in the assembler wrapper.\par\medskip
\noindent\textit{Editorial note.} Text is retained verbatim, including the author's own typos and placeholders.\par\medskip
\noindent\textit{Editorial note.} Queue index 3 of queue\_discriminant.json (pb=2725, stmt 2417-2423, proof 2425-2479).\par\medskip
\noindent\textit{Allowed standard prerequisites (editorial).}\begin{itemize}
\item Explicit prerequisite boundary. The cited inputs are prerequisites; the argument in the retained span is not replaced by any cited result.
\end{itemize}\medskip
\ifdefined\hypertarget\hypertarget{sourcepointer-1}{}\fi
\noindent\textit{Source reference.} Original lemma 04GG: sources/stacks-discriminant/pinned-original-source.tex, source lines 2479--2479. Cited by the retained statement or proof but not part of the collected spans; supplied only as a navigation pointer. The argument does not depend on its content.\par\medskip
\ifdefined\hypertarget\hypertarget{sourcepointer-2}{}\fi
\noindent\textit{Source reference.} Original lemma 03C3: sources/stacks-discriminant/pinned-original-source.tex, source lines 2479--2479. Cited by the retained statement or proof but not part of the collected spans; supplied only as a navigation pointer. The argument does not depend on its content.\par\medskip
\ifdefined\hypertarget\hypertarget{sourcepointer-3}{}\fi
\noindent\textit{Source reference.} Original lemma 0FKZ: sources/stacks-discriminant/pinned-original-source.tex, source lines 2326--2339. Cited by the retained statement or proof but not part of the collected spans; supplied only as a navigation pointer. The argument does not depend on its content.\par\medskip
\ifdefined\hypertarget\hypertarget{sourcepointer-4}{}\fi
\noindent\textit{Source reference.} Original lemma 0FK8: sources/stacks-discriminant/pinned-original-source.tex, source lines 2168--2185. Cited by the retained statement or proof but not part of the collected spans; supplied only as a navigation pointer. The argument does not depend on its content.\par\medskip
\ifdefined\hypertarget\hypertarget{sourcepointer-5}{}\fi
\noindent\textit{Source reference.} Original lemma 0FKY: sources/stacks-discriminant/pinned-original-source.tex, source lines 2282--2298. Cited by the retained statement or proof but not part of the collected spans; supplied only as a navigation pointer. The argument does not depend on its content.\par\medskip
\ifdefined\hypertarget\hypertarget{sourcepointer-6}{}\fi
\noindent\textit{Source reference.} Original lemma 0FL0: sources/stacks-discriminant/pinned-original-source.tex, source lines 2363--2368. Cited by the retained statement or proof but not part of the collected spans; supplied only as a navigation pointer. The argument does not depend on its content.\par\medskip
\ifdefined\hypertarget\hypertarget{sourcepointer-7}{}\fi
\noindent\textit{Source reference.} Original lemma 0FL1: sources/stacks-discriminant/pinned-original-source.tex, source lines 2384--2388. Cited by the retained statement or proof but not part of the collected spans; supplied only as a navigation pointer. The argument does not depend on its content.\par\medskip
\ifdefined\hypertarget\hypertarget{sourcepointer-8}{}\fi
\noindent\textit{Source reference.} Original lemma 0FK9: sources/stacks-discriminant/pinned-original-source.tex, source lines 2188--2195. Cited by the retained statement or proof but not part of the collected spans; supplied only as a navigation pointer. The argument does not depend on its content.\par\medskip
\ifdefined\hypertarget\hypertarget{sourcepointer-9}{}\fi
\noindent\textit{Source reference.} Original lemma 06LN: sources/stacks-discriminant/pinned-original-source.tex, source lines 2479--2479. Cited by the retained statement or proof but not part of the collected spans; supplied only as a navigation pointer. The argument does not depend on its content.\par\medskip
\ifdefined\hypertarget\hypertarget{sourcepointer-10}{}\fi
\noindent\textit{Source reference.} Original lemma 04HF: sources/stacks-discriminant/pinned-original-source.tex, source lines 2479--2479. Cited by the retained statement or proof but not part of the collected spans; supplied only as a navigation pointer. The argument does not depend on its content.\par\medskip
\ifdefined\hypertarget\hypertarget{sourcepointer-11}{}\fi
\noindent\textit{Source reference.} Original lemma 0476: sources/stacks-discriminant/pinned-original-source.tex, source lines 2479--2479. Cited by the retained statement or proof but not part of the collected spans; supplied only as a navigation pointer. The argument does not depend on its content.\par\medskip
\begin{lemma}
\label{lemma-universal-finite-syntomic-etale}
With notation as in Example \ref{example-universal-finite-syntomic}
consider the open subscheme $U'_d \subset X_d$ over which
$\pi_d$ is \'etale. Then $U'_d$ is a dense subset of the
open $U_d$ of Lemma \ref{lemma-universal-finite-syntomic}.
\end{lemma}

\begin{proof}
By exactly the same reasoning as in the proof of
Lemma \ref{lemma-universal-finite-syntomic}, using
Morphisms, Lemma \ref{morphisms-lemma-set-points-where-fibres-etale},
there is a maximal open $U'_d \subset X_d$ over which $\pi_d$ is
\'etale. Moreover, since an \'etale morphism is syntomic, we see
that $U'_d \subset U_d$. To finish the proof we have to show
that $U'_d \subset U_d$ is dense. Let $u : \Spec(k) \to U_d$ be a morphism
where $k$ is a field. Let $B = k \otimes_{A_d} B_d$ as in the
proof of Lemma \ref{lemma-universal-finite-syntomic-smooth}.
We will show there is a local domain $A'$ with residue field $k$
and a finite syntomic $A'$ algebra $B'$ with $B = k \otimes_{A'} B'$
whose generic fibre is \'etale. Exactly as in the previous paragraph
this will determine a morphism $\Spec(A') \to U_d$ which will map the
generic point into $U'_d$ and the closed point to $u$, thereby
finishing the proof.

\medskip\noindent
By Lemma \ref{lemma-syntomic-finite} part (4) we can choose a presentation
$B = k[x_1, \ldots, x_n]/(f_1, \ldots, f_n)$.
Let $d'$ be the maximum total degree of the polynomials $f_1, \ldots, f_n$.
Let $Y_{n, d'} \to X_{n, d'}$ be as in
Example \ref{example-universal-quasi-finite-syntomic}.
By construction there is a morphism $u' : \Spec(k) \to X_{n, d'}$
such that
$$
\Spec(B) \cong Y_{n, d'} \times_{X_{n, d'}, u'} \Spec(k)
$$
Denote $A = \mathcal{O}_{X_{n, d'}, u'}^h$ the henselization of the
local ring of $X_{n, d'}$ at the image of $u'$. Then we can write
$$
Y_{n, d'} \times_{X_{n, d'}} \Spec(A) = Z \amalg W
$$
with $Z \to \Spec(A)$ finite and $W \to \Spec(A)$ having empty
closed fibre, see
Algebra, Lemma \ref{algebra-lemma-characterize-henselian} part (13)
or the discussion in More on Morphisms, Section
\ref{more-morphisms-section-etale-localization}.
By Lemma \ref{lemma-universal-quasi-finite-syntomic-etale}
the local ring $A$ is regular (here we also use
More on Algebra, Lemma \ref{more-algebra-lemma-henselization-regular})
and the morphism $Z \to \Spec(A)$ is \'etale over the generic point of
$\Spec(A)$ (because it is mapped to the generic point of $X_{d, n'}$).
By construction $Z \times_{\Spec(A)} \Spec(k) \cong \Spec(B)$.
This proves what we want except that the map from
residue field of $A$ to $k$ may not be an isomorphism.
By Algebra, Lemma \ref{algebra-lemma-flat-local-given-residue-field}
there exists a flat local ring map $A \to A'$ such that the residue
field of $A'$ is $k$. If $A'$ isn't a domain, then we choose a
minimal prime $\mathfrak p \subset A'$ (which lies over the
unique minimal prime of $A$ by flatness) and we replace
$A'$ by $A'/\mathfrak p$. Set $B'$ equal to the unique $A'$-algebra
such that $Z \times_{\Spec(A)} \Spec(A') = \Spec(B')$.
This finishes the proof.
\end{proof}

\end{document}
